\chapter{Discussion}
\label{ch:discussion}

\section{What the three findings add up to}

Taken together:
\begin{itemize}
  \item \textbf{Detection (Chapter~\ref{ch:results-detection}).}  SOTA windowed models beat classical ones on same-dataset eTaPR; rankings change with metric choice; PA-F1 inflation is real.
  \item \textbf{Transfer (Chapter~\ref{ch:results-transfer}).}  Source-calibrated thresholds don't transfer; target-calibrated thresholds reveal a genuine classical/SOTA split; within-HAI LOPO shows feature-specific blindness.
  \item \textbf{Attribution (Chapter~\ref{ch:results-attribution}).}  Dense AE beats SOTA on localization; attention rollout adds nothing.
\end{itemize}

The single thesis-level story: \textbf{ICS anomaly detection is not a one-number benchmark.}  Different metrics, calibration regimes, and attribution tasks all reshuffle the ranking.  Any paper that picks one (model, metric, threshold) triple to report has under-constrained its own claim.  Reviewers should ask for the full grid; our contribution is a protocol for producing it.

\section{Likely reviewer objections}

\paragraph{``Your transfer projection is lossy.''}
True, by design (Chapter~\ref{ch:results-transfer}, \S1).  We accepted option~1 (per-variable type tagging) as the simplest defensible scheme.  Option~2 (learned projections) and option~3 (within-testbed-only) remained unused, and we report the ceiling-check against within-HAI LOPO: even in the most favourable within-testbed condition, the ceiling is $0.57$ pointwise F1 --- an order of magnitude better than cross-testbed but still far from same-dataset.  The cross-testbed numbers are not artifacts of a bad projection; they are upper-bounded by the problem itself.

\paragraph{``TranAD reproducibility gap.''}
PA-F1 $0.46$ vs.\ Tuli et al.~\cite{tuli2022tranad}'s $0.97$.  Sources enumerated in \S\ref{sec:tranad-repro}: dataset version, subsample, architecture capacity, threshold protocol.  Our TranAD still ranks \#1 on HAI PA-F1 in our setup --- the qualitative claim survives.  The thesis does not depend on matching published absolutes; it depends on comparing models under a consistent protocol.

\paragraph{``USAD collapsed on Morris --- selection bias?''}
USAD eTaPR $0.009$ on Morris (Chapter~\ref{ch:results-detection}) is not cherry-picked; it is what the model's scores produce.  The collapse is reported as a finding about USAD's assumptions on discrete data, not hidden.

\paragraph{``Only one seed for Phase 4.''}
Acknowledged; listed as a limitation.  Phase 1--2 multi-seed numbers showed eTaPR variance $\leq 0.02$ across seeds for the classical models; attribution seed-variance is similar in the multi-seed runs reported in Chapter~\ref{ch:results-attribution}.  The P1/P2/P3 ordering is robust to that magnitude of noise.

\paragraph{``Process-level attribution is weak.''}
Yes; HAI 21.03 does not publish per-sensor attack targets in machine-readable form.  Chapter~\ref{ch:results-attribution}, \S``Limitations'' lists this as the primary weakness and points at HAI 22.04 as a follow-up.

\section{Practical implications for deployment}

\begin{itemize}
  \item \textbf{Calibrate on target-normal.}  Any deployment should collect a clean-operation period from the target plant before activating the detector and calibrate the threshold there.  Source-trained percentile thresholds produce trivial classifiers (\S\ref{sec:finding-src-cal}).
  \item \textbf{Report three metrics.}  Readers should not trust single-metric rankings.  Our recommendation: pointwise F1 for per-timestep claims, PA-F1 for ``at least one timestep of the event was flagged'' (and flag the inflation per Kim et al.~\cite{kim2022towards}), eTaPR~\cite{hwang2022etapr} for event-level detection with coverage weighting.
  \item \textbf{Consider Dense AE for localization even if a SOTA model is the detector.}  A hybrid pipeline --- SOTA for flagging, Dense AE for attribution --- would outperform either alone on operator-facing triage.
\end{itemize}

\section{What we deliberately did not test}

\begin{itemize}
  \item New detector architectures.  The thesis is empirical, not architectural.
  \item SWaT~\cite{goh2017swat} / WADI~\cite{ahmed2017wadi}.  Adding more single-testbed datasets does not change the cross-testbed contribution.
  \item Real-time streaming performance.  Every number is offline batched.
  \item Federated or differentially-private training.  Out of scope.
\end{itemize}

\section{Future work}

\begin{enumerate}
  \item \textbf{HAI 22.04 per-sensor attribution.}  22.04's README contains per-attack target-sensor descriptions that can be parsed into a ground-truth mapping.  Running the Chapter~\ref{ch:results-attribution} eval against that would upgrade the localization claim from process-level to sensor-level.
  \item \textbf{Multi-seed attribution} across all six models to confirm or refute the LSTM-AE $\equiv$ USAD coincidence (Chapter~\ref{ch:results-attribution}, Finding 4.3).
  \item \textbf{Learned projection for schema alignment} (CLAUDE.md \S5.2 option 2) to test whether the cross-testbed ceiling moves under a richer alignment.  We expect modest improvement given the LOPO ceiling argument above, but the experiment is cheap to run.
  \item \textbf{PA-F1 replacement in the ICS AD literature.}  Our data reproduces the Kim et al.~\cite{kim2022towards} critique on ICS benchmarks; an editorial push to eTaPR as the default would be a small but useful community-level contribution.
\end{enumerate}
